% 連続体力学 連続体力学-3 問題1（018-grad） 解いた日 2026-10-28
%% kotae
\toi{1}{◎}{}
\[ \tens{D}=\begin{pmatrix}a&0&b\\0&0&0\\c&0&0\end{pmatrix} \]
\toi{2}{◎}{}
\[ \tens{E}=\begin{pmatrix}a&0&\frac{b+c}{2}\\0&0&0\\\frac{b+c}{2}&0&0\end{pmatrix} \]
\[ \tens{F}=\begin{pmatrix}0&0&\frac{b-c}{2}\\0&0&0\\\frac{c-b}{2}&0&0\end{pmatrix} \]
\toi{3}{◎}{}
\[ \Omega_{1}=0,\quad \Omega_{2}=b-c,\quad \Omega_{3}=0 \]
\begin{align*}
\vect{\nabla}\times\vect{r}&=\biggl(\pdv{}{x_{2}}cx_{1}-\pdv{}{x_{3}}0,\\
&\qquad\pdv{}{x_{3}}(ax_{1}+bx_{3})-\pdv{}{x_{1}}cx_{1},\\
&\qquad\pdv{}{x_{1}}0-\pdv{}{x_{2}}(ax_{1}+bx_{3})\biggr)\\
&=(0,\ b-c,\ 0)
\end{align*}
\toi{4}{◎}{}
\[ \tens{F}=0\quad\therefore\ b=c \]
\[ \tens{E}=0\quad a=0\text{ かつ }c=-b \]
$\tens{E}=0$ のとき
\[ \tens{F}=\begin{pmatrix}0&0&b\\0&0&0\\-b&0&0\end{pmatrix} \]
$x_{2}$ 軸　反時計　角度 $b$
%% kizuki
